% Generated draft: requires mathematical and source review.
\documentclass[11pt]{article}
\usepackage[margin=25mm]{geometry}
\usepackage{fontspec}
\setmainfont{Latin Modern Roman}
\usepackage{amsmath,amssymb,mathtools,unicode-math}
\setmathfont{Latin Modern Math}
\usepackage{enumitem,needspace,xurl,hyperref}
\setlength{\parindent}{0pt}
\setlength{\parskip}{5pt}
\setlength{\emergencystretch}{3em}
\newcommand{\N}{\mathbb{N}}
\newcommand{\Z}{\mathbb{Z}}
\newcommand{\Q}{\mathbb{Q}}
\newcommand{\R}{\mathbb{R}}
\newcommand{\C}{\mathbb{C}}
\newcommand{\F}{\mathbb{F}}
\newcommand{\A}{\mathbb{A}}
\newcommand{\PP}{\mathbb{P}}
\newcommand{\E}{\mathbb{E}}
\newcommand{\eps}{\varepsilon}
\newcommand{\dd}{\,\mathrm{d}}
\begin{document}
\section*{UnsolvedMath ID: 800938}
\textbf{Source problem code:} catalogue entry
\subsection{Definitions and mathematical statement}
For a Riemannian manifold $(M,g)$, write
\[
SM=\{(x,v)\in TM:g_x(v,v)=1\}
\]
for its unit tangent bundle. Whenever the unit-speed geodesics are defined for all real times, its geodesic flow is
\[
\phi_t(x,v)=\bigl(\gamma_{x,v}(t),\dot\gamma_{x,v}(t)\bigr),\qquad t\in\R,
\]
where $\gamma_{x,v}(0)=x$ and $\dot\gamma_{x,v}(0)=v$. The source's notation $S^M_i$ is interpreted as $SM_i$.

For a nonnegative integer $k$, the geodesic flows of $(M_1,g_1)$ and $(M_2,g_2)$ are \emph{$C^k$-conjugate} if there is a bijection $h:SM_1\to SM_2$ such that $h$ and $h^{-1}$ are $C^k$ and
\[
h\circ\phi_t^1=\phi_t^2\circ h\qquad\text{for every }t\in\R.
\]
In particular, a $C^0$ conjugacy is a homeomorphism satisfying this identity with the same time parameter.

\textbf{Conjecture 18 (Burns--Matveev, as attributed in the supplied record).} Let $(M_1,g_1)$ and $(M_2,g_2)$ be compact Riemannian manifolds with negative curvature, interpreted here as strictly negative sectional curvature:
\[
K_{g_i}(x,\sigma)<0
\]
for each $i\in\{1,2\}$, each $x\in M_i$, and every two-dimensional subspace $\sigma\subseteq T_xM_i$. Suppose their geodesic flows are defined for all $t\in\R$ and are $C^0$-conjugate. Then there exists a smooth diffeomorphism $F:M_1\to M_2$ such that
\[
F^*g_2=g_1.
\]
Thus the two Riemannian manifolds are isometric; the statement does not require the given conjugacy $h$ to equal the map induced by $F$.

\textbf{Unresolved source qualifications.} The supplied excerpt does not explicitly specify a boundary convention, a dimension restriction, or connectedness. Its all-time flow identity requires globally defined geodesic flows, which is made explicit above. No boundary dynamics are supplied. The sectional-curvature interpretation is the usual interpretation in this setting, but the excerpt itself says only ``negative curvature.'' The supplied registry labels this a longstanding open conjecture; that status is not independently verified here.
\subsection{Short English statement}
The conjecture attributed to Burns and Matveev says that compact negatively curved Riemannian manifolds are isometric whenever a homeomorphism between their unit tangent bundles intertwines their geodesic flows for every real time, preserving the time parameter. Boundary, dimension, and curvature conventions require human review against the source.
\subsection{Sources}
\begin{itemize}
\item Local UnsolvedMath registry: permanent ID $800938$, source problem code \texttt{catalogue entry}, title \emph{Open Problems and Questions about Geodesics --- Conjecture 18}. The supplied record attributes the problem to Burns and Matveev, gives proposed year 2019, and labels its status \texttt{open}; these are supplied registry assertions.
\item Burns, Matveev, \emph{Open Problems and Questions about Geodesics} (2019, as identified in the supplied source list), Conjecture 18. Supplied source URL: \url{https://arxiv.org/abs/1308.5417}.
\item Supplied provenance: accessed 2026-10-16; extraction method \texttt{arxiv\_latex}; upstream revision \texttt{d371ee917ac9647fbb06cd6fc529deec6f3dda82}. The access date and source description are supplied provenance, not a claim of independent source consultation or literature verification.
\end{itemize}
\end{document}
